\documentclass[11pt,a4paper]{article}

\usepackage[margin=1in]{geometry}
\usepackage[T1]{fontenc}
\usepackage[utf8]{inputenc}
\usepackage{lmodern}
\usepackage{microtype}
\usepackage{graphicx}
\usepackage{float}
\usepackage{booktabs}
\usepackage{xcolor}
\usepackage{listings}
\usepackage[hidelinks]{hyperref}

\input{glyphtounicode}
\pdfgentounicode=1

\definecolor{codebg}{RGB}{248,248,248}
\definecolor{rulegray}{RGB}{210,210,210}

\lstset{
  basicstyle=\ttfamily\footnotesize,
  breaklines=true,
  breakatwhitespace=false,
  columns=fullflexible,
  keepspaces=true,
  frame=single,
  rulecolor=\color{rulegray},
  backgroundcolor=\color{codebg},
  xleftmargin=0.2in,
  xrightmargin=0.2in
}

\title{Cryptography Engineering Project 2 Report}
\author{Group 9}
\date{}

\begin{document}
\maketitle

\section{Task 1: Single Hash vs. Merkle Tree}

\subsection{Experimental Results}

\textbf{Test Environment Configuration}
\begin{itemize}
  \item Dataset size: 128 MB
  \item Number of blocks: $n = 2^7 = 128$
\end{itemize}

\textbf{Single Hash (All-or-Nothing Approach)}
\begin{itemize}
  \item Resulting SHA-256:
\end{itemize}
\begin{lstlisting}
0c05acf58bb36284a9e73f0222f2fd0621954bb8eb5c16ee50a45341b5cd4594
\end{lstlisting}
\begin{itemize}
  \item Execution time: 0.088621 seconds
\end{itemize}

\textbf{Merkle Tree Construction}
\begin{itemize}
  \item Merkle root:
\end{itemize}
\begin{lstlisting}
3cd74d8d7171cce47fb900a1aca739b74286b6f84e5da9490f54559dbecc8668
\end{lstlisting}
\begin{itemize}
  \item Execution time: 0.137790 seconds
  \item Internal node memory overhead: 8255 bytes
\end{itemize}

\begin{figure}[H]
  \centering
  \includegraphics[width=0.92\linewidth]{report_assets/image1.png}
  \caption{Task 1 console output.}
\end{figure}

\subsection{Discussion and Findings}

\textbf{Execution Time Comparison.}
The Single Hash method completed in approximately 0.088 seconds, while the
Merkle Tree construction took 0.137 seconds. The Single Hash approach is
inherently faster because it sequentially processes the continuous 128 MB data
stream in a single pass. In contrast, the Merkle Tree incurs a slight time
penalty. This delay is expected, as the algorithm must slice the file into 128
distinct blocks, compute 128 individual leaf hashes, and then execute iterative
loops to concatenate and hash the internal nodes upward to the root.

\textbf{Memory Overhead Analysis.}
To fulfill the requirement of analyzing the memory overhead of the tree's
internal nodes, the measured overhead is 8255 bytes. Mathematically, a binary
Merkle Tree with $n = 128$ leaf nodes requires $n - 1 = 127$ internal parent
nodes. Storing 127 SHA-256 hash objects in Python aligns with this roughly 8 KB
footprint.

\section{Task 2: Efficient Error Localization}

\subsection{Implementation and Experiment Results}

\textbf{Top-down search algorithm.}
In order to find the corrupted block, the implementation compares the Merkle
tree from the original correct data with the Merkle tree from the corrupted
data. Starting from the root node, each traversal level needs only one child
comparison because the experiment assumes that only one block is corrupted. If
the left child hash does not match the trusted tree, the corruption is located
in the left subtree, so the algorithm traverses down the left path. If the left
child hash matches, the left subtree is intact and the corruption must be in the
right subtree, so the algorithm traverses down the right path. This process
continues until the affected leaf node is reached.

\begin{figure}[H]
  \centering
  \includegraphics[width=0.98\linewidth]{report_assets/image2.png}
  \caption{Task 2 console output.}
\end{figure}

\subsection{Discussion}

In this experiment, the 128 MB file is divided into $n = 128$ blocks, which are
used to construct a complete binary tree. The height $H$ of this tree is
$\log_2(128) = 7$.

\textbf{Traversal Strategy.}
During the top-down traversal, because there is only a single bit error in the
entire dataset, it is sufficient to check only one child node, specifically the
left child, at each level to make a routing decision.

\textbf{Comparison Count Proof.}
Each hash comparison eliminates half of the remaining search space. Traversing
from the root at level 7 down to a leaf at level 0 requires descending exactly
7 levels. Because the implementation executes exactly one hash comparison per
level, the total number of comparisons is exactly 7. This satisfies the
theoretical constraint of $\log_2(n)$ comparisons.

\section{Task 3: Efficient Node Replacement}

\subsection{Experimental Results}

\begin{figure}[H]
  \centering
  \includegraphics[width=0.95\linewidth]{report_assets/image3.png}
  \caption{Task 3 console output.}
\end{figure}

\subsection{Discussion and Findings}

Task 3 evaluates whether the Merkle root can be updated efficiently after
replacing one 1 MB block. Instead of rebuilding the entire Merkle Tree,
\texttt{replace\_node()} updates only the affected path from the modified leaf
to the root.

\textbf{Correctness Verification.}
Both update methods produced the same updated Merkle root.

\begin{lstlisting}
c392a826e9a141f3f4dfef8baf007620a1b0d16e6fd888da5a3556b2d863039e
\end{lstlisting}

Since both methods produced the same root, the path update implementation is
correct.

\textbf{Hash Computation Comparison.}
The full reconstruction method recomputes the entire Merkle Tree, requiring 128
leaf hashes and 127 internal hashes, for a total of 255 hash computations. In
contrast, \texttt{replace\_node()} only recomputes the modified leaf and the 7
internal nodes on the path to the root, requiring only 8 hash computations.
Therefore, the path update reduces the number of hash computations by:

\[
\frac{255}{8} = 31.88
\]

\textbf{Runtime Comparison.}
The measured path update time was 0.000989 seconds, while the full
reconstruction time was 0.078458 seconds.

\[
\frac{0.078458}{0.000989} = 79.33
\]

Therefore, the path update was approximately 79.33 times faster in this
experiment. This demonstrates the main advantage of Merkle Trees for dynamic
data. When only one block changes, only the path from the modified leaf to the
root needs to be updated, giving an $O(\log n)$ update cost instead of
recomputing the entire file.

\section{Task 4: Advanced Self-Healing}

\subsection{Implementation and Experiment Results}

Task 4 integrates XOR parity recovery with Merkle Tree error detection. The
input files are the corrupted 128 MB dataset, the trusted Merkle Tree generated
from the original data, and the parity blocks generated from the original data.
The clean original file is not used during repair.

The parity file stores one parity block for every pair of data blocks:

\[
P_i = D_{2i} \oplus D_{2i+1}
\]

After \texttt{locate\_error()} identifies the corrupted block, the implementation
selects the corresponding parity block and the sibling data block. If block
$D_j$ is corrupted, its sibling block $D_s$ is still correct because the
experiment contains only one bit flip in one block. The original block can
therefore be reconstructed as:

\[
D_j = P_{\lfloor j/2 \rfloor} \oplus D_s
\]

The recovered block is then passed to \texttt{replace\_node()}, which repairs
only the affected leaf-to-root path in the corrupted Merkle Tree. The final
verification compares the repaired root against the trusted root.

\begin{lstlisting}
========================================
 Task 4: Advanced Self-Healing
========================================
1. Loading corrupted data, trusted Merkle Tree, and parity blocks...
2. Building Merkle Tree from corrupted data...
3. Locating corrupted block using trusted Merkle Tree...
4. Recovering corrupted block with XOR parity and sibling block...
5. Repairing Merkle Tree with replace_node()...

[Output]
Trusted Merkle Root        : 3cd74d8d7171cce47fb900a1aca739b74286b6f84e5da9490f54559dbecc8668
Corrupted File Root        : bb2b791c16712525185c86f30f39bc878365f4612423a106ad0f729f5003242b
Corrupted Block Index      : 42
Comparison Count           : 7
Parity Block Index Used    : 21
Sibling Block Index        : 43
Repaired Merkle Root       : 3cd74d8d7171cce47fb900a1aca739b74286b6f84e5da9490f54559dbecc8668
Verification Result        : PASS
========================================
\end{lstlisting}

\subsection{Discussion and Findings}

The corrupted block was identified as block 42 using only the corrupted file and
the trusted Merkle Tree. Since $n = 128$, the tree height is
$\log_2(128) = 7$, and the reported comparison count is exactly 7. This confirms
that the detection step keeps the required $O(\log n)$ localization behavior.

The parity index used for recovery was 21 because block 42 belongs to the block
pair $(42, 43)$, and $\lfloor 42/2 \rfloor = 21$. Since block 42 is the even
member of the pair, block 43 is the corresponding sibling block. The repair
process computes:

\[
D_{42} = P_{21} \oplus D_{43}
\]

The repaired Merkle root is identical to the trusted Merkle root:

\begin{lstlisting}
3cd74d8d7171cce47fb900a1aca739b74286b6f84e5da9490f54559dbecc8668
\end{lstlisting}

Therefore, the corrupted block can be both located and recovered without access
to the original clean 128 MB file. The Merkle Tree provides efficient
localization, while XOR parity provides the missing data needed for correction.
Together, they implement a self-healing mechanism for a single corrupted block.

\section{Summary Table}

\begin{table}[H]
\centering
\small
\begin{tabular}{p{0.18\linewidth}p{0.23\linewidth}p{0.25\linewidth}p{0.22\linewidth}}
\toprule
Operation & Single Hash & Merkle Tree & Improvement Factor \\
\midrule
Detecting error & Scans 128 MB & Checks the trusted root against the corrupted root & Root mismatch gives immediate detection after root computation \\
Localizing error & Impossible & 7 comparisons & $\log_2(128)$ search \\
Updating 1 block & Computes 128 MB & Computes 1 MB block hash plus 7 internal hashes & 31.88x fewer hashes \\
Proof size & No local proof & $7 \times 32 = 224$ bytes & Logarithmic proof size \\
\bottomrule
\end{tabular}
\caption{Comparison summary for $n = 2^7$ blocks.}
\end{table}

\end{document}
